\section{Class Diagram Specification}
\label{appendix:class-diagram}

This appendix expands the three domain packages summarised in Section~\ref{fig:class_diagram}. For each package it lists the entities, the relationships that carry design intent, and --- where a reader might reasonably expect a different shape --- the reason the model does not take it.

% ────────────────────────────────────────────────────────────
\subsection{Identity \& Access}

\subsubsection{Roles are assignments, not subclasses}

A natural first reading of an education platform is that \texttt{Student} and \texttt{Teacher} are subclasses of an abstract \texttt{User}. The system deliberately does not do this. There is one concrete \texttt{User} entity holding identity and authentication state; every capability is derived at request time from role assignments.

Three properties of the domain make inheritance the wrong tool:

\begin{itemize}
    \item \textbf{One person holds several roles at once.} A doctoral teaching assistant is an instructor on one course and an enrolled student on another. Under inheritance that person needs two rows, two identities, and two sets of credentials.
    \item \textbf{A role is scoped, not global.} ``Instructor'' is never true in the abstract --- it is true \emph{of a particular course}. \texttt{user\_role\_assignments} carries a \texttt{scope\_kind} (\texttt{global}, \texttt{organization}, \texttt{org\_unit}, \texttt{course}) plus the matching foreign key, so the same role means different things in different places. A subclass has no room to hold that scope.
    \item \textbf{Roles change without the person changing.} Granting and revoking is an insert and a soft delete; under inheritance it would be a change of class, which no ORM models comfortably and which would invalidate every existing reference to the row.
\end{itemize}

\begin{table}[H]
\centering
\renewcommand{\arraystretch}{1.25}
\small
\begin{tabular}{|>{\raggedright\arraybackslash}m{0.26\textwidth}|>{\raggedright\arraybackslash}m{0.66\textwidth}|}
\hline
\textbf{Entity} & \textbf{Role in the model} \\ \hline
\texttt{User} & Identity and authentication state. One concrete class; no subclasses. \\ \hline
\texttt{UserProfile} & Display name, avatar and contact details, separated so that identity can exist before a profile is completed. \\ \hline
\texttt{Role} & A named bundle of permissions (\texttt{code}, e.g.\ \texttt{course\_instructor}). \\ \hline
\texttt{Permission} & The atomic capability a guard checks, e.g.\ \texttt{course.update}. \\ \hline
\texttt{UserRoleAssignment} & Binds a user to a role \emph{at a scope}. The join carrying \texttt{scope\_kind} and the scope's foreign key. \\ \hline
\texttt{UserPermissionGrant} & A direct grant, for the exceptional case that does not justify a role. \\ \hline
\texttt{AuthSession} & An issued session; the unit that MFA and logout act on. \\ \hline
\texttt{AuthEvent} & The typed authentication and access-control audit trail (append-only). \\ \hline
\end{tabular}
\caption{Identity \& Access entities.}
\label{tab:class_identity}
\end{table}

\noindent A permission check therefore resolves as: \emph{user} $\rightarrow$ live assignments whose scope contains the target resource $\rightarrow$ the union of those roles' permission codes $\rightarrow$ does that set contain the required code?

% ────────────────────────────────────────────────────────────
\subsection{Course \& Ingestion}

\subsubsection{File bytes are not stored in the database}

Material content lives in object storage (S3-compatible), never in a database column. A \texttt{LearningMaterialVersion} holds \texttt{storage\_object\_id}, a foreign key to a \texttt{StorageObject} row that records the bucket and key; the bytes themselves are fetched from, and streamed to, the object store.

The reason is operational rather than stylistic. Lecture material is measured in tens or hundreds of megabytes per file; holding it in Postgres would push every backup, replica and point-in-time restore to the same scale, and every read of a row would carry the risk of dragging the payload with it. Object storage also gives presigned, expiring URLs, so a client downloads directly from the store without the API service proxying the bytes.

\begin{table}[H]
\centering
\renewcommand{\arraystretch}{1.25}
\small
\begin{tabular}{|>{\raggedright\arraybackslash}m{0.30\textwidth}|>{\raggedright\arraybackslash}m{0.62\textwidth}|}
\hline
\textbf{Entity} & \textbf{Role in the model} \\ \hline
\texttt{Course} & The container a learner enrols in; owns modules and materials. \\ \hline
\texttt{Module} $\rightarrow$ \texttt{Lesson} & The ordered curriculum tree beneath a course. \\ \hline
\texttt{LearningMaterial} & A logical document attached to a lesson (``Week 3 slides''), independent of which upload is current. \\ \hline
\texttt{LearningMaterialVersion} & One uploaded revision, carrying \texttt{storage\_object\_id} and the processing status of the ingestion pipeline. \\ \hline
\texttt{StorageObject} & Bucket and key of the stored bytes. The boundary between the database and the object store. \\ \hline
\texttt{DocumentChunk} & A retrievable passage of a version, with its text, metadata, content hash, and a 3072-dimensional embedding. \\ \hline
\end{tabular}
\caption{Course \& Ingestion entities.}
\label{tab:class_ingestion}
\end{table}

\noindent Versioning is what allows a material to be re-uploaded without invalidating the work derived from it: the old version's chunks and embeddings remain queryable until the new version finishes processing, so retrieval never falls into a gap mid-ingestion.

% ────────────────────────────────────────────────────────────
\subsection{Assessment \& Interview}

\begin{table}[H]
\centering
\renewcommand{\arraystretch}{1.25}
\small
\begin{tabular}{|>{\raggedright\arraybackslash}m{0.30\textwidth}|>{\raggedright\arraybackslash}m{0.62\textwidth}|}
\hline
\textbf{Entity} & \textbf{Role in the model} \\ \hline
\texttt{Quiz} $\rightarrow$ \texttt{QuizQuestion} & The assessment definition and its items. \\ \hline
\texttt{QuizAttempt} & One learner's run at a quiz; the unit grading and integrity events attach to. \\ \hline
\texttt{QuizAttemptAnswer} & The per-question response within an attempt. \\ \hline
\texttt{InterviewConfig} & The teacher-authored definition of a mock interview. \\ \hline
\texttt{InterviewSession} & One learner's run, holding the conversation and its evaluation. \\ \hline
\texttt{GapReport} & The aggregated outcome of a session: which learning outcomes the answers did and did not evidence. \\ \hline
\texttt{AssessmentIntegrityEvent} & Browser integrity signals recorded during an attempt or session; append-only. \\ \hline
\end{tabular}
\caption{Assessment \& Interview entities.}
\label{tab:class_assessment}
\end{table}

\noindent The two assessment modes meet at the \texttt{GapReport}: quiz attempts establish whether a learner can recall and apply the material, the interview session probes whether they can explain it, and the report reconciles both against the course's declared learning outcomes.
